\begin{textsample}{POS Dim 1 – unprofiled\_gpt – Score 6.00 – unprofiled\_gpt/t000124\_gpt.txt}  \label{ex:f1_pos_025}
well I know they ’re going to challenge me, that ’s a given.
Of course they will.
They ’re my friends, and they ’re going to be my biggest critics.
So I think the only \textbf{way} I can deal with that properly is to be really \textbf{clear} about why I ’m \textbf{changing} direction and what I ’m actually trying to do.
I ’ll have to spell out that I ’m in a kind of transitional phase, not just randomly flitting about.
I ’m moving towards creating something \textbf{different}, even if from the outside it might just look a bit vague or experimental right now.
For me it *is* new, because I ’m learning new things all the time – stuff I didn't know before, things I ’m finding out and learning about as I go.
It feels new because I ’m stretching myself.
A big part of it is learning from people who are already doing what I want to do.
I want to make it \textbf{clear} to them that I ’m not just dreaming this up in isolation ; I ’m watching and listening to people who are actually out there doing it, and I ’m trying to pick up as much as I can from them.
That ’s how I ’m thinking about it : learning from others, then gradually shaping my \textbf{own} version.
I ’ll probably explain that in entrepreneurship there are common themes, but everyone has their \textbf{own} path.
There ’s no single route you ’re supposed to follow.
So yes, there are shared \textbf{ideas} and patterns, but what I ’m doing is finding my \textbf{own} \textbf{way} through that, in a \textbf{way} that makes \textbf{sense} for me.
Even if my friends don't fully get my perspective, or they can't quite see the “ bigger picture ” I ’ve got in my head, I still think there ’s value in learning from them too.
They might not understand why I ’m doing this, but they ’ll still have experiences and insights that are useful.
So my response to their criticism will be to calmly lay out my \textbf{reasons}, show that I ’m in a deliberate transition towards something \textbf{different}, and keep stressing that this is about learning – from people ahead of me, and from them as well – even if they don't entirely see what I see yet.

% matched lemmas: change, clear, different, idea, own, reason, sense, way
\end{textsample}
